\documentclass[12pt]{article}
\usepackage{amsmath, amsthm, amsfonts, mathtools, xcolor, caption}

\usepackage{minimalJ}

\numberwithin{equation}{section}

% Requires Biber
\usepackage[style=numeric,sorting=none,giveninits=true]{biblatex}
\addbibresource{bibliography.bib}

\usepackage{etoolbox}
\usepackage[margin=1in]{geometry}

\DeclareMathAlphabet{\mymathbb}{U}{BOONDOX-ds}{m}{n}

%\usepackage[colorlinks, allcolors=blue, hypertexnames=false]{hyperref}
%\usepackage[noabbrev, nameinlink]{cleveref}

\usepackage{comment}

%%%%%%%%%%%%%%%%%%%%%%%%%%%%%%%%%%%%%%%%%
%       Theorem Environments            %
%%%%%%%%%%%%%%%%%%%%%%%%%%%%%%%%%%%%%%%%%

\newtheorem{theorem}{Theorem}[section]
\newtheorem{lemma}[theorem]{Lemma}
\newtheorem{corollary}[theorem]{Corollary}
\newtheorem{proposition}[theorem]{Proposition}
\theoremstyle{definition}
\newtheorem{remark}{Remark}
\newtheorem{definition}{Definition}
\newtheorem{notation}{Notation}
\newtheorem{example}{Example}
\newtheorem{conjecture}[theorem]{Conjecture}
\newtheorem{openproblem}[theorem]{Open Problem}

%%%%%%%%%%%%%%%%%%%%%%%%%%%%%%%%%%%%%%%%%
%               New Commands            %
%%%%%%%%%%%%%%%%%%%%%%%%%%%%%%%%%%%%%%%%%

\newcommand{\ev}[1]{( #1 )} 
\newcommand{\ew}[1]{\left( #1 \right)} 

\newcommand{\N}{\mathbb{N}}
\newcommand{\Z}{\mathbb{Z}}
\newcommand{\C}{\mathbb{C}}
\newcommand{\Q}{\mathbb{Q}}
\newcommand{\R}{\mathbb{R}}



%%%%%%%%%%%%%%%%%%%%%%%%%%%%%%%%%%%%%%%%%
%               Document                %
%%%%%%%%%%%%%%%%%%%%%%%%%%%%%%%%%%%%%%%%%

\let\underbrace\LaTeXunderbrace
\let\overbrace\LaTeXoverbrace

\begin{document}

\title{Leanified Theorems}

%\date{\today}
%\maketitle

%%%%%%%%%%%%%%%%%%%%%%%%%%%%%%%%%%%%%%%%%%%%%%%%%%%%%%%%%%%%%%%%%%%%%%%%%%%%%%%%%%%%%%%%%%%%%%%%%%%%%%%%%%%%%
\section*{Challenges}
\refstepcounter{section}
%%%%%%%%%%%%%%%%%%%%%%%%%%%%%%%%%%%%%%%%%%%%%%%%%%%%%%%%%%%%%%%%%%%%%%%%%%%%%%%%%%%%%%%%%%%%%%%%%%%%%%%%%%%%%

\begin{enumerate}
\item \textbf{The Hadwiger Conjecture. }

%Let $G= \ev{V,E}$ be a finite simple graph.

\begin{definition}[Graph Minor] 
Given simple graphs $X$ and $G$, we say that $X$ is a \emph{minor} of $G$ if there is a function assigning to each vertex $x$ of $X$ a nonempty vertex set $V_x\subseteq V(G)$ such that 
\begin{enumerate}
\item the subgraph $G[V_x]$ induced by $G$ on $V_x$ is connected for each $x$;
\item for different vertices  $x,y\in V\ev{X}$, the sets $V_x$ and $V_y$ are disjoint;
\item for any two adjacent vertices $x,y\in V\ev{X}$, there are adjacent vertices $a\in V_x$ and $b\in V_y$ of $G$. 
\end{enumerate}
\end{definition}

COMMENT FROM NATHAN: maybe we can skip the definitions of minor and chromatic number

\begin{definition}[Hadwiger Number]
Given a graph $G$, the \textit{Hadwiger number} of $G$ is the largest integer $n$ such that $K_n$ is a minor of $G$; here $K_n$ denotes the complete graph on $n$ vertices. 
We denote the Hadwiger number by $h\ev{G}$.
\end{definition}

\begin{definition}[Chromatic number]
Given a graph $G$  and $n\in \mathbb{N}$, 
a \textit{proper coloring} of $G$ with $n$ colours is a function $f:[n] \rightarrow V\ev{G}$ so that adjacent vertices receive different $f$-values. The \textit{chromatic number} of $G$ is the minimum number $n$ such that there is a proper coloring of $G$ with $n$ colours. 
\end{definition}


\begin{conjecture}[Hadwiger's Conjecture
\mbox{\cite[Conjecture 1]{steiner}}]\label{hadw}
Let $r\in \mathbb{N}$.
For every finite simple graph $G$, we have:
$$
r \leq \chi\ev{G} \ \Longrightarrow \  r \leq h\ev{G} 
$$
\end{conjecture}

While Conjecture \ref{hadw} is not difficult for every small values of $r$, already for $r=4$ it is equivalent to the 4-colour theorem. 
Robertson, Seymour and Thomas proved the above for $r= 5$ \cite{RobertsonSeymourThomas1993Hadwiger}. 

\newpage
\item \textbf{Excluded Minors of $\mathrm{GF}(p^m)$-representable matroids.}

\begin{definition}[Matroid]
A \emph{matroid} is a pair $M=(E,\mathcal I)$ consisting of a finite set $E$ and a collection $\mathcal I$ of subsets of $E$, called the \emph{independent sets}, such that:
\begin{enumerate}
\item $\emptyset \in \mathcal I$;
\item if $I\in \mathcal I$ and $J\subseteq I$, then $J\in \mathcal I$;
\item if $I,J\in \mathcal I$ and $|I|<|J|$, then there is some $x\in J\setminus I$ such that $I\cup \{x\}\in \mathcal I$.
\end{enumerate}
The maximal independent sets of $M$ are called the \emph{bases} of $M$.
\end{definition}

\begin{definition}[$\mathrm{GF}(p^m)$-representable matroid]
Let $p$ be a prime number. Let $m$ be a positive integer. Let $\mathrm{GF}(p^m)$ denote the finite field with $p^m$ elements. A matroid $M$ is \emph{$\mathrm{GF}(p^m)$-representable} if there is a matrix $A$ over $\mathrm{GF}(p^m)$ such that the independent sets of $M$ are exactly the sets of columns of $A$ which are linearly independent over $\mathrm{GF}(p^m)$.
\end{definition}

\begin{definition}[Matroid minor]
Let $M$ be a matroid with ground set $E$. The \textit{dual} of $M$, denoted $M^*$, is the matroid with the same ground set, where the bases are the complements of the bases of $M$. Given a subset $D \subseteq E$, the \emph{deletion} of $D$ from $M$, denoted $M\setminus D$, is the restriction of $M$ to $E\setminus D$. Given a subset $C\subseteq E$, the \emph{contraction} of $M$ onto $E-C$ is the matroid $M/C := (M^* \setminus C)^*$, with ground set $E\setminus C$. 

Given matroids $M,N$, then $N$ is a \emph{minor} of $M$ if there exist sets $C$ and $D$ such that $N = M/C \setminus D$.
\end{definition}

\begin{definition}[Excluded minor]
Let $\mathcal C$ be a class of matroids closed under minors. A matroid $M$ is an \emph{excluded minor} for $\mathcal C$ if $M\notin \mathcal C$, but every proper minor of $M$ belongs to $\mathcal C$.
\end{definition}

The class of $\mathrm{GF}(p^m)$-representable matroids is closed under taking minors. Thus, it is natural to ask for minimal matroids which fail to be representable over $\mathrm{GF}(p^m)$.

\begin{openproblem}
Let $p$ be a prime number. Let $m$ be a positive integer. List the minimal excluded minors for the class of $\mathrm{GF}(p^m)$-representable matroids and prove that the given list is complete up to isomorphism.
\end{openproblem}

For $p^m=2$, the excluded minor list is known: a matroid is representable over $\mathrm{GF}(2)$ if and only if it has no minor isomorphic to $U_{2,4}$ \cite{Tutte1958Homotopy}. For $p^m=3$, there is also a finite explicit list: a matroid is representable over $\mathrm{GF}(3)$ if and only if it has no minor isomorphic to $U_{2,5}$, $U_{3,5}$, $F_7$, or $F_7^*$ \cite{Bixby1979Reid,Seymour1979GF3}. In general, Rota's excluded-minor conjecture, proven by Geelen, Gerards and Whittle, says that for every finite field $\mathrm{GF}(p^m)$ there are only finitely many excluded minors for the class of $\mathrm{GF}(p^m)$-representable matroids \cite{Rota1971Combinatorial, GeelenGerardsWhittle2014Solving}. 

We have the following parametrized version of the open problem.
\begin{openproblem}
Let $r$ be a positive integer. If there is a prime number $p$ and a positive integer $m$ such that $r=p^m$, then list the minimal excluded minors for the class of $\mathrm{GF}(p^m)$-representable matroids and prove that the list is complete up to isomorphism.
\end{openproblem}

In the formal statement, the list is given as a finite set $L$ of matroids: the members of $L$ are pairwise non-isomorphic excluded minors, and every finite matroid that is not $\mathrm{GF}(p^m)$-representable has a minor isomorphic to a member of $L$. This rules out indirect solutions that prove the existence of a complete list without exhibiting one.

TODO: Check references. 

\begin{comment}
\newpage
\item \textbf{Excluded Minors of $GF(r)$-representable matroids.}
%\cite{}

\begin{openproblem}
Let $r$ be a natural number which is a prime power. List the minimal excluded minors for the class of GF$(p^m)$-representable matroids and prove that the given list is complete up to isomorphism.
\end{openproblem}
\end{comment}

\newpage
\item {\bf Ramsey Numbers.}
\begin{definition}[Ramsey Number]
Given $t\in \mathbb{N}$, the \emph{(diagonal) Ramsey number} \( R(t) \) is the smallest $n\in \mathbb{N}\cup \{\infty\}$ such that every simple graph \( G \) on \( n \) vertices contains either a clique of size \( t \) or an independent set of size \( t \).
\end{definition}
By Ramsey's theorem, $R(t)$ is always an integer, which is upper bounded by $4^t$ \cite{ramsey1930formal, ramseygraham1990ramsey}.
Erd\H{o}s proved that random graphs $G(n,p)$ give a lower bound of the form $\sqrt{2}^t$ for $t$ sufficently large \cite{erdos1947some}. An important problem  that can be traced back to Ramsey's original paper \cite{ramsey1930formal} is to get good upper and lower bounds for the Ramsey number.

\begin{openproblem}[Scaled Ramsey problem]
Let $r\in \mathbb{N}$.
Provide decimal numbers $d_1$ and $d_2$ with $|d_2-d_1|\leq (4-\sqrt{2}) \cdot 0.9^r$ so that for all sufficiently large $t\in \mathbb{N}$
\[
d_1^t \leq R(t) \leq d_2^t.
\]
\end{openproblem}

The scaled Ramsey problem for $r=1$ follows by combining the above mentioned works of Ramsey and Erd\H{o}s.
The problem was solved for $r=2$ by Campos, Griffiths, Morris  and Sahasrabudhe
\cite{campos2026exponential} combined with the better bound from the follow-up paper by 
Gupta, Ndiaye, Norin and Wei \cite{gupta2024optimizing}.

\newpage
\item \textbf{Sidorenko Conjecture for Half-Graphs.} 
% \cite{ }

\begin{definition}[Graph homomorphism density]
Let $G$ and $H$ be simple graphs. 
%\phantom{a}
\begin{enumerate}
\item A function $f: V\ev{H} \rightarrow V\ev{G}$ is a \textit{graph homomorphism} from $H$ to $G$ if for each edge $uv \in E\ev{H}$, we have $f\ev{u}f\ev{v} \in E\ev{G}$. 

\item Let $\mathrm{hom}\ev{H,G}$ be the number of graph homomorphisms $H \rightarrow G$. 

\item The \textit{homomorphism density} of $H$ in $G$ is the fraction
$$
t_{H}\ev{G} := \dfrac{\mathrm{hom}\ev{H,G}}{|V\ev{G}|^{|V\ev{H}|}}
$$
\end{enumerate}
\end{definition}

\begin{conjecture}[Sidorenko's Conjecture \mbox{\cite{sidorenko}}]
For every bipartite graph $H$ and for any graph $G$, we have:
$$
t_H\ev{G} \geq t_{K_2}\ev{G}^{|E\ev{H}|}
$$
\end{conjecture}


We are interested in the special case of the conjecture, where $H = H_r$ is a half-graph with $2r$ vertices, as follows. As far as we are aware, the conjecture has not been asked anywhere for half-graphs, and the only reason we do it for half-graphs here is that we believe that they are somewhat generic, and it will help us formulate an easier scalable version. (TODO really ask some experts in extremal combinatorics here)


\begin{definition}[Half-graph]
Let $A_r = \{ a_1,\ldots, a_r \}$ and $B_r = \{ b_1,\ldots, b_r \} $ be disjoint sets of size $r$. 
A \emph{half-graph} is the bipartite graph on the vertex set $A_r\cup B_r$ with the following edges present: $E = \{ a_i b_j : i \leq j \}$
(Note that these are about \lq half\rq\ of the possible edges between $A_r$ and $B_r$, hence the name). 
\end{definition}

\begin{conjecture}[Scaled weakening of Sidorenko's Conjecture]
Let $r \in \N$. For any graph $G$, we have:
$$
t_{H_r}\ev{G} \geq t_{K_2}\ev{G}^{(r^2+r)/2}
$$
\end{conjecture}

TODO find out for which half-graphs it is known.

\newpage
\item \textbf{Erdős–Hajnal Conjecture.}


\begin{definition}[Cliques and Independent sets]
Given a simple graph $G$, a set of vertices of $G$ is a \textit{clique} if its vertices are pairwise adjacent. It is an \textit{independent set} if its vertices are pairwise non-adjacent. 
\end{definition}

\begin{definition}[$H$-free graphs]
Given graphs $G$ and $H$, we say that $G$ is \textit{$H$-free} if it does not contain an induced subgraph isomorphic to $H$. 
\end{definition}

\begin{conjecture}[Erdős–Hajnal Conjecture \mbox{\cite{erdos_hajnal}}]
For each graph $H$, there is a constant $c_H > 0$ such that every $H$-free graph $G$ with $n$ vertices has a clique or independent set of size at least $n^{c_H}$. 
\end{conjecture}



Given $r\in \mathbb{N}$, we denote the path with $r$ vertices by $P_r$.
We are interested in the following parametrized version of the conjecture, where $H=P_r$ is a path. 
\begin{conjecture}[Erdős–Hajnal Conjecture for paths]
Let $r\in \mathbb{N}$. There is some $c_r > 0$ such that every $P_r$-free graph $G$ with $n$ vertices has a clique or independent set of size $n^{c_r}$.
\end{conjecture}

The conjecture has been proved for $r=5$ by Nguyen, Tung, Scott and Seymour in \cite{nguyen2026induced}



\newpage
\item \textbf{Better-Quasi-Ordering of Finite Graphs under Minor Relation.}
\begin{definition}[Cumulative Hierarchy over $Q$]
Given a set $Q$, we denote by $\mathcal P^*$ the set of nonempty subsets of $Q$. We define for each ordinal $\alpha$ a set $V^*_{\alpha}(Q)$ recursively as follows:
\begin{itemize}
\item $V^*_0(Q) := Q$.
\item $V^*_{\beta + 1} := \cal{P}^* V^*_{\beta}$
\item $V^*_{\lambda}(Q) := \dot \bigcup_{\alpha < \lambda} V^*_{\alpha}(Q)$ for any limit ordinal $\lambda$.
\end{itemize}

The classs given by the union of all these sets is denoted denoted $V^*(Q)$. 
\end{definition}

\begin{definition}[$\alpha$-well-quasiordering]
Given a quasi-order $(Q,\le)$, we define a relation $\le^*$ on $V^*(Q)$ in terms of a 2-player game between Player I and Player II. A position in this game is a pair $(X,Y)$ with $X$ and $Y$ in $V_*(Q)$. If $X$ and $Y$ are both in $Q$ then this position is a win for Player II if $X \le Y$ and a win for Player I otherwise. If not, then Player I names a set $X'$, which must be $X$ itself if $X \in Q$ and otherwise must be an element of $X$. Then Player II also names a set $Y'$, which must be $Y$ itself if $Y \in Q$ and otherwise must be an element of $Y$. After this, the new position is $(X', Y')$. We define $X \le^* Y$ to mean that Player II has a winning strategy in this game with starting position $(X,Y)$.

We say that $(Q, \le)$ is {\em $\alpha$-well-quasiordered} if the restriction of this game to $V^*_{\alpha}(Q)$ is well-quasiordered.
\end{definition}

Note the $(Q, \le)$ is better-quasi-ordered if and only if it is $\omega_1$-well-quasi-ordered.

\begin{conjecture}
Let $r$ be an ordinal number. If $r$ is even, of the form $\alpha . 2$, then finite planar graphs are $\alpha$-well-quasi-ordered with respect to the minor relation. If $r$ is odd, of the form $\alpha . 2 + 1$, then finite graphs are $\alpha$-well-quasi-ordered.
\end{conjecture}

\newpage
\item {\bf Rota's Basis Conjecture.}

The following conjecture in its original form \cite[Conjecture~4]{huang1994relations} is stated in terms of matroids, but even in the special case where the matroid is a vector matroid -- that is, the bases come from a finite-dimensional vector space -- the conjecture is widely open.



\begin{conjecture}[Rota's basis conjecture]
Let $M$ be a finite matroid with a family $(B_1,..,B_n)$ of disjoint bases.
There is a family of disjoint bases $(C_1,..C_n)$ of $M$ so that $B_i\cap C_j$ is a singleton for all $i,j\in [n]$.
\end{conjecture}

Here we are interested in the following scaled weakening of the original conjecture.

\begin{conjecture}[Scaled weakening of Rota's basis conjecture]
Let $r\in \mathbb{N}$.
Let $M$ be a finite matroid with a family $(B_1,..,B_n)$ of disjoint bases. 
There is $n_0\in \mathbb{N}$ so that for all $n\geq n_0$ and all $\epsilon\in \mathbb{R}_{>0}$, there is a family of disjoint bases $(C_1,..C_{\lceil (1-1/r-\epsilon) n\rceil})$ of $M$ so that $B_i\cap C_j$ is a singleton for all $i,j\in [n]$.
\end{conjecture}

Buci\'c, Kwan, Pokrovskiy and Sudakov proved the scaled weakening for $r=2$ \cite{bucic2020halfway}.

\newpage
\item \textbf{Ryser's Hypergraph Conjecture}.


\begin{definition}[$r$-partite]
Given $r\in \mathbb{N}$, an $r$-uniform hypergraph $\mathcal{H}$ is \emph{$r$-partite} if there is a partition
$$
V\ev{\mathcal{H}} = \bigcup_{i=1}^{r} P_i
$$
of the vertex set of $\mathcal{H}$ such that $|P_i \cap e| = 1$, for each $e\in E\ev{\mathcal{H}}$ and $i\in [r]$. 
\end{definition}


\begin{definition}[Cover number]
Given a hypergraph $\mathcal{H}$, a
set $C$ of vertices of $\mathcal{H}$ is a \textit{vertex cover} if it contains a vertex of every edge. The smallest cardinality of a vertex cover of $\mathcal{H}$ is called the \textit{cover number} of $\mathcal{H}$. It is denoted by $\tau(\mathcal{H})$.
\end{definition}

\begin{definition}[Matching number]
The \textit{matching number} of a hypergraph $\mathcal{H}$ is the largest cardinality of a set of pairwise disjoint edges of $\mathcal{H}$. It is denoted by $\nu\ev{\mathcal{H}}$.
\end{definition}

\begin{conjecture}[Ryser's Conjecture \mbox{\cite{bishnoi}}]
Let $r\in \mathbb{N}$.
Every $r$-partite $r$-uniform hypergraph $\mathcal{H}$ satisfies:
$$
\tau\ev{\mathcal{H}} \leq \ev{r-1} \cdot \nu\ev{\mathcal{H}}
$$

\end{conjecture}

For $r=2$, it is a classical theorem of K\H{o}nig, and for $r=3$ it was proved by Aharoni \cite{AharoniRyserTripartite} using the topological method. 

\newpage
\item \textbf{Union-Closed Sets Conjecture.} 

\begin{definition}[Union-closed set]
A set $\mathcal{F}$ of sets is \textit{union-closed} if for any two sets in $\mathcal{F}$, their union is also in $\mathcal{F}$. 
For a nonempty finite set $\Fcal$, the \emph{density} of an element $x$ with respect to $\Fcal$ is defined
to be 
\[
\frac{|\{ F \in \mathcal{F} : x \in F \}|}{|\Fcal|}
\]
\end{definition}

\begin{conjecture}[Union-closed sets conjecture \mbox{\cite{nagel2023}}]
For every finite union-closed set $\mathcal{F}\notin\{\emptyset, \{\emptyset\}\}$, there exists an element $x$ of density at least one half.
\end{conjecture}

We are interested in the following parametrized version of the conjecture.
\begin{conjecture}[Scaled union-closed sets conjecture]
Let $r\in \mathbb{N}$.
For every finite union-closed set $\mathcal{F}\notin\{\emptyset, \{\emptyset\}\}$, there exists an element $x$ of density at least $\frac{1}{2}-\frac{1}{r+2}$.
\end{conjecture}

For $r=0$ the bound equals $0$ and the statement is elementary, while as $r\to\infty$ the bound approaches $\frac{1}{2}$ and the statement recovers the full conjecture in the limit. The best published lower bound on the density ($\approx 0.382$, obtained via the entropy method) settles the cases $r\leq 6$; the conjecture is open for all $r\geq 7$.






\newpage
\item \textbf{The Unfriendly Partition Conjecture} (Cowan and Emerson, see for example \cite{aharoni1990unfriendly})

\begin{definition}[Partial partition]
A {\em partial partition} of the vertices of a graph $G$ is a partial function $V(G) \to \{0,1\}$. We call it {\em finite} if its domain is finite. We call it a {\em partition} if it is a total function.
\end{definition}

\begin{definition}[Unfriendly]
Let $v$ be a vertex of a graph $G$. A partial partition $f \colon V(G) \to \{0,1\}$ is {\em unfriendly at $v$} if $v$ has at least as many neighbours mapped by $f$ to the opposite partition class as it has neighbours mapped to the same partition class. Formally this means
$$|N(v) \cap f^{-1}[1 - f(v)]| \ge |N(v) \cap f^{-1}[f(v)]|.$$
A partition is {\em unfriendly} if it is unfriendly at every vertex.
\end{definition}

\begin{conjecture}[Unfriendly partition Conjecture]
Every countable graph has an unfriendly partition.
\end{conjecture}

In order to pose a scaled version of this conjecture, we introduce the following game:

\begin{definition}[The $r$-partitioning game.]
Let $r$ be an ordinal number and let $G$ be a graph. The {\em $r$-partitioning game on $G$} is played between two players, Challenger and Partitioner, as follows. A position in the game is a pair $(f, \alpha)$ consisting of a partial partition $f$ of $G$ and an ordinal number $\alpha$. The initial position is $(\emptyset, r)$. In each turn, beginning in position $(f, \alpha)$, Challenger must name a vertex $v$ of $G$ and an ordinal $\beta < \alpha$ and Partitioner must then name a finite partial partition $g$ extending $f$ which is defined at $v$ and unfriendly at $v$. After this, the new position is $(g, \beta)$. If at any point some player is unable to move, the other player wins.
\end{definition}

Note that since there are no infinite decreasing sequences of ordinals this game must end after finitely many turns. 

The scaled conjecture is as follows:

\begin{conjecture}[Scaled Unfriendly Partition Conjecture]
Let $r$ be an ordinal. For any countable graph $G$, Partitioner has a winning strategy in the $r$-partitioning game on $G$.
\end{conjecture}
\end{enumerate}

\newpage

%%%%%%%%%%%%%%%%%%%%%%%%%%%%%%%%%%%%%%%%%%%%%%%%%%
\printbibliography 
%%%%%%%%%%%%%%%%%%%%%%%%%%%%%%%%%%%%%%%%%%%%%%%%%%
\end{document}